\subsection{特征多项式的补充}

设 L 是一个交换环，M 是一个有限秩 m 的自由 L-模。若 u 是 M 的一个自同态，r 是一个自然数，则我们用 \(c_{r}(u)\) 表示自同态 \(\wedge^{r}(u)\)（自由 L-模 \(\wedge^{r}(\mathrm{M})\) 的）的迹。特别地，我们有

\[
c _ {0} (u) = 1, \quad c _ {1} (u) = \mathrm{Tr} (u), \quad c _ {m} (u) = \det (u),\tag{1}
\]

并且 \(c_{r}(u) = 0\)（对 r \textgreater{} m）。由 III, §8, No.~4, 第 527 页的命题 7，映射 \(u \mapsto \det(u)\)（从 End(M) 到 L 的）是 m 次齐次多项式映射（IV, §5, No.~9, 第 55 页）。更一般地，对满足 \(0 \leqslant r \leqslant m\) 的每个整数 r，映射 \(c_{r}\)（从 End(M) 到 L 的）是 r 次齐次多项式映射；这由 III, §8, No.~5, 第 529 页的命题 10 得出。

设 \(u\) 是 \(M\) 的一个自同态，\(\overline{u}\) 是 \(L[X]\)-模 \(M[X] = M \otimes_{L} L[X]\) 的自同态（由 \(u\) 经标量扩张得到，II, §5, No.~1, 第 277 页）。回顾（III, §8, No.~11, 第 541 页，定义 3 与 (50)），\(u\) 的特征多项式是行列式 \(\chi_u(X)\)（\(L[X]\)-自同态 \(X - \overline{u}\)（\(M[X]\) 的）的），并且我们有关系式

\[
\chi_ {u} (\mathrm{X}) = \sum_ {r = 0} ^ {m} (- 1) ^ {r} c _ {r} (u) \mathrm{X} ^ {m - r}.\tag{2}
\]

命题 1. —— 设 L 是一个交换环，M 是一个有限秩 \(m \geqslant 1\) 的自由 L-模，u 是 M 的一个自同态。存在 M 的唯一一个自同态 \(\widetilde{u}\) 满足关系式

\[
\widetilde {u} (x) \wedge w = x \wedge \wedge^ {m - 1} (u) (w)\tag{3}
\]

（对 \(x\in \mathrm{M}\) 与 \(w\in \bigwedge^{m - 1}(\mathrm{M})\)）。此外，我们有关系式

\[
u \circ \widetilde {u} = \widetilde {u} \circ u = \det (u) _ {\mathrm{M}},\tag{4}
\]

\[
\det (\widetilde {u}) = \det (u) ^ {m - 1},\tag{5}
\]

\[
\widetilde {u} = \sum_ {r = 0} ^ {m - 1} (- 1) ^ {r} c _ {m - 1 - r} (u) u ^ {r}.\tag{6}
\]

引理 1. —— 设 p 是一个满足 \(0 \leqslant p \leqslant m\) 的整数。对 \(\wedge^{p}(\mathrm{M})\) 中的任何 w，设 \(h_{p}(w)\) 是线性映射 \(w' \mapsto w \wedge w'\)（从 \(\wedge^{m-p}(\mathrm{M})\) 到 \(\wedge^{m}(\mathrm{M})\) 的）。线性映射 \(h_{p}: w \mapsto h_{p}(w)\)（从 \(\wedge^{p}(\mathrm{M})\) 到 \(\operatorname{Hom}_{\mathrm{L}}(\wedge^{m-p}(\mathrm{M}), \wedge^{m}(\mathrm{M}))\) 的）是一个同构。

设 \((e_{i})_{i\in\mathrm{I}}\) 是 M 的一个基；我们给集合 I 赋予一个全序。对 I 的任何子集 J，令 \(e_{\mathrm{J}}=e_{i_{1}}\wedge\cdots\wedge e_{i_{r}}\)，其中 \((i_{1},\ldots,i_{r})\) 是 J 的元素按递增次序排成的序列。L-模 \(\wedge^{m-p}(\mathrm{M})\) 以元素 \(e_{\mathrm{S}}\) 为基，其中 S 取遍 I 的具有 m-p 个元素的子集所成之集；\(\wedge^{m}(\mathrm{M})\) 以 \(\{e_{\mathrm{I}}\}\) 为基。因此，\(\mathrm{Hom}_{\mathrm{L}}(\wedge^{m-p}(\mathrm{M}),\wedge^{m}(\mathrm{M}))\) 有一个由线性映射 \(e_{\mathrm{J}}^{*}\) 组成的基，它们由公式

\[
e _ {\mathrm{J}} ^ {*} (e _ {\mathrm{S}}) = \left\{ \begin{array}{l l} e _ {\mathrm{I}} & \text {if } I = J\cup S, \\ 0 & \text {otherwise}, \end{array} \right.\tag{7}
\]

刻画，其中 J 取遍 I 的具有 p 个元素的子集所成之集。由 III, §7, No.~8, 第 519 页的公式 (20)，对 I 的每个具有 p 个元素的子集 J，我们有 \(h_{p}(e_{\mathrm{J}}) \in \{e_{\mathrm{J}}^{*}, -e_{\mathrm{J}}^{*}\}\)；由于元素 \(e_{J}\) 构成 \(\Lambda^{p}(\mathrm{M})\) 的一个基，线性映射 \(h_{p}\) 是双射。

现在来证明命题 1。设 \(u\) 与 \(\widetilde{u}\) 是 M 的自同态。关系式 (3) 等价于

\[
h _ {1} \circ \widetilde {u} = \operatorname{Hom} (\wedge^ {m - 1} (u) \cdot 1 _ {\bigwedge^ {m} (\mathrm{M})}) \circ h _ {1};\tag{8}
\]

由引理 1，映射 \(h_{1}\) 是从 M 到 \(\operatorname{Hom}_{\mathrm{L}}(\wedge^{m-1}(\mathrm{M}),\wedge^{m}(\mathrm{M}))\) 的一个同构。因此，对 M 的每个自同态 u，存在 M 的唯一一个满足关系式 (3) 的自同态 \(\widetilde{u}\)。

设 \(x_{1},\ldots,x_{m}\) 是 M 的元素。在 (3) 中把 x 换成 \(u(x_{1})\)，把 w 换成 \(x_{2}\wedge\cdots\wedge x_{m}\)；我们得到

\[
\widetilde {u} (u (x _ {1})) \wedge x _ {2} \wedge \dots \wedge x _ {m} = u (x _ {1}) \wedge \dots \wedge u (x _ {m}) = \det (u) x _ {1} \wedge \dots \wedge x _ {m}.
\]

因此 \(h_1(\widetilde{u} \circ u(x_1)) = h_1(\det(u)x_1)\)，由引理 1 这给出关系式 \(\widetilde{u} \circ u = \det(u)_\mathrm{M}\)。

我们用 U 表示自同态 \(X - \overline{u}\)（L{[}X{]}-模 M{[}X{]} 的）（VIII, 第 335 页）。把上面应用于 U，存在一个自同态 \(\widetilde{U}\)（L{[}X{]}-模 M{[}X{]} 的），它满足关系式

\[
\widetilde {\mathrm{U}} (x _ {1}) \wedge x _ {2} \wedge \dots \wedge x _ {m} = x _ {1} \wedge (\mathrm{X} x _ {2} - u (x _ {2})) \wedge \dots \wedge (\mathrm{X} x _ {m} - u (x _ {m}))\tag{9}
\]

（对 M 中的 \(x_{1},\ldots ,x_{m}\)）以及

\[
\widetilde {\mathrm{U}} \circ \mathrm{U} = \det (\mathrm{X} - \overline {{u}}) _ {\mathrm{M} [ \mathrm{X} ]}.\tag{10}
\]

我们把 \(\widetilde{\mathbf{U}}\) 看作 \(\operatorname{End}(\mathbf{M})[\mathbf{X}]\) 的一个元素（VIII, 第 9 页）；由公式 (9) 与引理 1，它的次数 \(\leqslant m - 1\)，因此我们可以把它写成

\[
\widetilde {\mathrm{U}} = \sum_ {r = 0} ^ {m - 1} (- 1) ^ {r} u _ {r} X ^ {m - 1 - r},\tag{11}
\]

其中 \(u_{r}\) 是 M 的自同态。由公式 (2)，关系式 (10) 给出等式

\[
\left(\sum_ {r = 0} ^ {m - 1} (- 1) ^ {r} u _ {r} \mathrm{X} ^ {m - 1 - r}\right) (\mathrm{X} - u) = \sum_ {r = 0} ^ {m} (- 1) ^ {r} c _ {r} (u) \mathrm{X} ^ {m - r}\tag{12}
\]

（在环 \(\operatorname{End}(M)[X]\) 中）。通过让两边的 X 的单项式的系数相等，我们得到关系式

\[
u _ {r} + u _ {r - 1} \circ u = c _ {r} (u) _ {\mathrm{M}} \qquad \mathrm{for} 1 \leqslant r \leqslant m - 1\tag{13}
\]

以及

\[
u _ {0} = c _ {0} (u) _ {\mathrm{M}}, \quad u _ {m - 1} \circ u = c _ {m} (u) _ {\mathrm{M}}.\tag{14}
\]

由此我们推出

\[
u _ {m - 1} = \sum_ {r = 0} ^ {m - 1} (- 1) ^ {r} c _ {m - 1 - r} (u) u ^ {r}.\tag{15}
\]

现在，当确定常数项时，等式 (9) 与 (11) 给出 \(u_{m-1} = \widetilde{u}\)；由此得出公式 (6)。

特别地，\(\widetilde{u}\) 属于 \(\operatorname{End}(\mathbf{M})\) 的由 \(u\) 生成的子代数，从而与 \(u\) 交换。我们已经建立了关系式 \(\widetilde{u} \circ u = \det(u)_{\mathrm{M}}\)；由此得出公式 (4)。

最后，设 \(x_{1},\ldots,x_{m}\) 是 M 的元素。在公式 (3) 中我们把 x 换成 \(x_{1}\)，把 w 换成 \(\widetilde{u}(x_{2})\wedge\cdots\wedge\widetilde{u}(x_{m})\)。这给出

\[
\begin{array}{c} \widetilde {u} (x _ {1}) \wedge \widetilde {u} (x _ {2}) \wedge \dots \wedge \widetilde {u} (x _ {m}) = x _ {1} \wedge u \circ (\widetilde {u} (x _ {2})) \wedge \dots \wedge u \circ (\widetilde {u} (x _ {m})) \\ = \det (u) ^ {m - 1} x _ {1} \wedge x _ {2} \wedge \dots \wedge x _ {m} \end{array}
\]

从而得出公式 (5)。

注. —— 1) M 的自同态 \(\widetilde{u}\) 与我们在 III, §8, No.~6, 第 532 页所称的 \(u\) 的余转置重合。

\begin{enumerate}
\def\labelenumi{\arabic{enumi})}
\setcounter{enumi}{1}
\item
  由公式 (1)、(2)、(4) 与 (6)，我们推出关系式 \(\chi_u(u) = 0\)，从而得到凯莱–哈密顿定理的另一个证明（III, §8, No.~11, 第 541 页）。
\item
  由于映射 \(c_{r}\)（从 End(M) 到 L 的）是一个次数为 r 的齐次多项式映射（对 \(0 \leqslant r \leqslant m - 1\)），由公式 (6) 推出，映射 \(u \mapsto \widetilde{u}\)（从 End(M) 到 End(M) 的）是一个次数为 m - 1 的齐次多项式映射。
\end{enumerate}

设 B 是环 L 上的一个代数；假定 B 是一个秩为 \(m \geqslant 1\) 的自由 L-模，并把 L 与 B 的子环 \(L \cdot 1\) 等同起来。设 b 是 B 的一个元素。我们把上面的结果应用于 L-模 B 的自同态 \(\gamma(b) : x \mapsto bx\)。令 \(\gamma_{r}(b) = c_{r}(\gamma(b))\)（对 \(0 \leqslant r \leqslant m\)）；特别地，我们有 \(\gamma_{m}(b) = \mathrm{N}_{\mathrm{B/L}}(b)\)（III, §9, No.~3, 第 543 页）。b 的特征多项式（出处同上）可以写成

\[
\mathrm{Pc} _ {\mathrm{B/L}} (b; \mathrm{X}) = \sum_ {r = 0} ^ {m} (- 1) ^ {r} \gamma_ {r} (b) \mathrm{X} ^ {m - r}.\tag{16}
\]

由于映射 \(\gamma\)（从 B 到 \(\operatorname{End}_{\mathrm{L}}(\mathrm{B})\) 的）是 L-线性的，映射 \(\gamma_{r}\)（从 B 到 L 的）是一个次数为 r 的齐次多项式映射。特别地，映射 \(b \mapsto \mathrm{N}_{\mathrm{B}/\mathrm{L}}(b)\)（从 B 到 L 的）是一个次数为 m 的齐次多项式映射。

对 B 的任何元素 \(b\)，令

\[
\widetilde {b} = \sum_ {r = 0} ^ {m - 1} (- 1) ^ {r} \gamma_ {m - 1 - r} (b) b ^ {r}.\tag{17}
\]

由命题 1，线性映射 \(\gamma (\widetilde{b})\)（从 B 到 B 的）是映射 \(\gamma (b)\) 的余转置；由此我们推出

\[
b \widetilde {b} = \widetilde {b} b = \mathrm{N} _ {\mathrm{B/L}} (b).\tag{18}
\]

此外，映射 \(b \mapsto \widetilde{b}\)（从 B 到 B 的）是一个次数为 \(m - 1\) 的齐次多项式映射。

